\documentclass[11pt, letterpaper]{article}

% --- Idioma y codificación ---
\usepackage[spanish]{babel}
\usepackage[utf8]{inputenc}
\usepackage[T1]{fontenc}
\usepackage{lmodern}

% --- Matemáticas ---
\usepackage{amsmath, amssymb, amsthm}

% --- Formato ---
\usepackage[margin=2.5cm]{geometry}
\usepackage{parskip}
\usepackage{booktabs}
\usepackage{hyperref}
\hypersetup{colorlinks=true, linkcolor=blue!70!black, urlcolor=blue!70!black}

% --- Entornos de teoremas ---
\theoremstyle{definition}
\newtheorem{definition}{Definición}[section]
\theoremstyle{plain}
\newtheorem{theorem}{Teorema}[section]
\newtheorem{lemma}[theorem]{Lema}
\newtheorem{corollary}[theorem]{Corolario}
\theoremstyle{remark}
\newtheorem{remark}{Observación}[section]

\title{%
  T03 --- Demostración formal:\\[4pt]
  Cota de altura de árboles rojinegros
}
\author{Curso Linux--C--Rust --- B15 Estructuras de datos en C}
\date{}

\begin{document}
\maketitle
\tableofcontents

% ===================================================================
\section{Propiedades rojinegras}
% ===================================================================

Un árbol rojinegro es un BST donde cada nodo tiene color rojo o negro
y se cumplen:

\begin{enumerate}
  \item[\textbf{P1.}] Cada nodo es rojo o negro.
  \item[\textbf{P2.}] La raíz es negra.
  \item[\textbf{P3.}] Cada hoja (NULL) es negra.
  \item[\textbf{P4.}] Si un nodo es rojo, ambos hijos son negros.
  \item[\textbf{P5.}] Para cada nodo, todos los caminos desde él hasta
    las hojas NULL contienen el mismo número de nodos negros.
\end{enumerate}

\begin{definition}[Black-height]
El black-height de un nodo $x$, denotado $\operatorname{bh}(x)$, es el
número de nodos negros en cualquier camino desde $x$ (sin incluir a $x$)
hasta una hoja NULL. P5 garantiza que este valor está bien definido.
\end{definition}

% ===================================================================
\section{Demostración 1: subárbol con
  \texorpdfstring{$\geq 2^{\operatorname{bh}(x)} - 1$}{>= 2\^bh(x) - 1}
  nodos internos}
% ===================================================================

\begin{theorem}\label{thm:subtree}
Sea $x$ un nodo en un árbol rojinegro. El subárbol enraizado en $x$
contiene al menos $2^{\operatorname{bh}(x)} - 1$ nodos internos.
\end{theorem}

\begin{proof}
Por inducción sobre la altura de $x$.

\textbf{Base (altura 0).} Si $x$ es una hoja NULL:
$\operatorname{bh}(x) = 0$ y el subárbol tiene 0 nodos internos:
\[
  2^{\operatorname{bh}(x)} - 1 = 2^0 - 1 = 0 \quad \checkmark
\]

\textbf{Paso inductivo.} Sea $x$ un nodo interno con hijos $l$ y $r$,
ambos con altura menor que la de $x$. Por hipótesis inductiva:
\begin{itemize}
  \item Subárbol de $l$: $\geq 2^{\operatorname{bh}(l)} - 1$ nodos.
  \item Subárbol de $r$: $\geq 2^{\operatorname{bh}(r)} - 1$ nodos.
\end{itemize}

Determinamos $\operatorname{bh}$ de los hijos en función de
$\operatorname{bh}(x)$:
\begin{itemize}
  \item Si el hijo es \textbf{negro}: $\operatorname{bh}(\text{hijo})
    = \operatorname{bh}(x) - 1$.
  \item Si el hijo es \textbf{rojo}: $\operatorname{bh}(\text{hijo})
    = \operatorname{bh}(x)$.
\end{itemize}
En ambos casos: $\operatorname{bh}(\text{hijo}) \geq \operatorname{bh}(x) - 1$.

Nodos internos en el subárbol de $x$:
\begin{align*}
  &\geq (2^{\operatorname{bh}(l)} - 1)
      + (2^{\operatorname{bh}(r)} - 1) + 1 \\
  &\geq (2^{\operatorname{bh}(x)-1} - 1)
      + (2^{\operatorname{bh}(x)-1} - 1) + 1 \\
  &= 2 \cdot 2^{\operatorname{bh}(x)-1} - 1
   = 2^{\operatorname{bh}(x)} - 1 \qedhere
\end{align*}
\end{proof}

% ===================================================================
\section{Demostración 2:
  \texorpdfstring{$\operatorname{bh}(r) \geq h/2$}{bh(r) >= h/2}}
% ===================================================================

\begin{theorem}\label{thm:bh-half}
Sea $T$ un árbol rojinegro con altura $h$ y raíz $r$. Entonces
$\operatorname{bh}(r) \geq h/2$.
\end{theorem}

\begin{proof}
En un camino raíz $\to$ NULL de longitud $h$, sean $b$ los nodos
negros bajo la raíz y $r_c$ los rojos. Entonces
$b + r_c = h$ y $\operatorname{bh}(r) = b$.

Por P4 (no dos rojos consecutivos), cada nodo rojo debe tener un hijo
negro debajo, lo que implica $r_c \leq b$. Entonces:
\[
  h = b + r_c \leq b + b = 2b
  \quad\Longrightarrow\quad
  b \geq h/2
  \quad\Longrightarrow\quad
  \operatorname{bh}(r) \geq h/2 \qedhere
\]
\end{proof}

% ===================================================================
\section{Demostración 3:
  \texorpdfstring{$h \leq 2 \cdot \log_2(n+1)$}{h <= 2 log2(n+1)}}
% ===================================================================

\begin{theorem}\label{thm:height}
Todo árbol rojinegro con $n$ nodos internos tiene altura
$h \leq 2 \cdot \log_2(n+1)$.
\end{theorem}

\begin{proof}
Sea $T$ un rojinegro con $n$ nodos internos, altura $h$, raíz $r$.

\textbf{Paso 1.} Por el Teorema~\ref{thm:subtree}:
$n \geq 2^{\operatorname{bh}(r)} - 1$.

\textbf{Paso 2.} Por el Teorema~\ref{thm:bh-half}:
$\operatorname{bh}(r) \geq h/2$.

\textbf{Paso 3.} Combinando:
\[
  n \geq 2^{\operatorname{bh}(r)} - 1 \geq 2^{h/2} - 1
\]
\[
  n + 1 \geq 2^{h/2}
  \quad\Longrightarrow\quad
  \log_2(n+1) \geq h/2
  \quad\Longrightarrow\quad
  h \leq 2\log_2(n+1) \qedhere
\]
\end{proof}

% ===================================================================
\section{Comparación formal con AVL}
% ===================================================================

\begin{corollary}
Para $n$ nodos, las cotas de altura son:
\end{corollary}

\begin{table}[h]
\centering
\begin{tabular}{lcc}
\toprule
\textbf{Estructura} & \textbf{Cota de altura} & \textbf{Constante} \\
\midrule
Árbol perfecto     & $\log_2(n{+}1)$          & 1.0  \\
AVL                & $1.44 \cdot \log_2(n{+}2)$ & 1.44 \\
Rojinegro          & $2 \cdot \log_2(n{+}1)$    & 2.0  \\
BST sin balance    & $n - 1$                   & ---  \\
\bottomrule
\end{tabular}
\end{table}

\begin{remark}[La cota es ajustada (tight)]
Un rojinegro con $\operatorname{bh} = k$ puede tener camino más largo
$2k$ (alternancia rojo-negro) y más corto $k$ (todo negro). La razón
es exactamente $2{:}1$, confirmando que el factor~2 no se puede mejorar.
\end{remark}

% ===================================================================
\section{Demostración 4: operaciones rojinegras son
  \texorpdfstring{$O(\log n)$}{O(log n)}}
% ===================================================================

\begin{corollary}\label{cor:ops}
Búsqueda, inserción y eliminación en un árbol rojinegro con $n$ nodos
tienen complejidad $O(\log n)$ en el peor caso.
\end{corollary}

\begin{proof}
\emph{Búsqueda}: camino raíz $\to$ hoja con $O(1)$ por nivel.
Costo $= O(h) = O(\log n)$. \checkmark

\emph{Inserción}: inserción BST estándar $O(h)$, más correcciones
que suben $O(h)$ niveles con $O(1)$ por nivel y \textbf{a lo sumo
2~rotaciones}. Total: $O(\log n)$. \checkmark

\emph{Eliminación}: eliminación BST $O(h)$, más correcciones con
\textbf{a lo sumo 3~rotaciones} y $O(h)$ recoloreos.
Total: $O(\log n)$. \checkmark
\end{proof}

\begin{remark}
Aunque la cota rojinegra ($2\log_2 n$) es peor que la AVL
($1.44\log_2 n$), la inserción y eliminación requieren menos
rotaciones en el peor caso ($2$--$3$ vs $O(\log n)$ en AVL). En la
práctica, los rojinegros suelen ser más rápidos en escenarios con
muchas modificaciones.
\end{remark}

% ===================================================================
\section{Resumen de resultados}
% ===================================================================

\begin{table}[h]
\centering
\begin{tabular}{llp{5.5cm}}
\toprule
\textbf{Resultado} & \textbf{Técnica} & \textbf{Clave} \\
\midrule
Subárbol $\geq 2^{\operatorname{bh}(x)}{-}1$ nodos
  & Inducción sobre altura
  & $\operatorname{bh}(\text{hijo}) \geq \operatorname{bh}(x){-}1$ \\
$\operatorname{bh}(r) \geq h/2$
  & P4 (no dos rojos seguidos)
  & $r_c \leq b \Rightarrow h \leq 2b$ \\
$h \leq 2\log_2(n{+}1)$
  & Combinar anteriores
  & $n \geq 2^{h/2} - 1$ \\
Cota tight
  & Alternancia rojo-negro
  & Razón camino largo/corto $= 2{:}1$ \\
\bottomrule
\end{tabular}
\end{table}

\end{document}
